% !TeX program = pdflatex
% -------------------------------------------------------------------------
%  crowbar_thyristor_selftrigger.tex
%  Crowbar thyristor for the 3 kV / 3600 A static-switch test bench:
%  device selection and passive self-triggering gate circuit.
%  Conventions: English, thebibliography (no biblatex), no babel, no lmodern.
% -------------------------------------------------------------------------
\documentclass[11pt,a4paper]{article}

\usepackage[T1]{fontenc}
\usepackage[utf8]{inputenc}
\usepackage[margin=25mm]{geometry}
\usepackage{amsmath,amssymb}
\usepackage{siunitx}
\usepackage{booktabs}
\usepackage{graphicx}
\usepackage{caption}
\usepackage[european]{circuitikz}
\usepackage{fancyhdr}
\usepackage[hidelinks]{hyperref}

% ---------------- document identification / revision ----------------------
\newcommand{\doctitle}{Crowbar Thyristor for the 3~kV / 3600~A Static-Switch Test Bench}
\newcommand{\docsubtitle}{Device Selection and Passive Self-Triggering Gate Circuit}
\newcommand{\docrev}{00}
\newcommand{\docdate}{2026-09-09}
\newcommand{\docid}{crowbar\_thyristor\_selftrigger}

\pagestyle{fancy}
\fancyhf{}
\lhead{\small \docid}
\rhead{\small Rev.~\docrev}
\lfoot{\small \docdate}
\cfoot{\small \thepage}
\renewcommand{\headrulewidth}{0.4pt}

\sisetup{per-mode=symbol, range-phrase=--, range-units=single}

\title{\textbf{\doctitle}\\[6pt]\large \docsubtitle\\[10pt]\normalsize Rev.~\docrev\ --- \docdate}
\author{}
\date{}

\begin{document}
\maketitle
\thispagestyle{fancy}

\begin{center}
\begin{tabular}{@{}lll@{}}
\toprule
\textbf{Rev.} & \textbf{Date} & \textbf{Description} \\
\midrule
00 & 2026-09-09 & First issue. Device selection, stress analysis, passive gate circuit. \\
\bottomrule
\end{tabular}
\end{center}

\begin{abstract}
A thyristor crowbar with a series dump inductor and resistor extracts the energy
of a \SI{2}{mH}, \SI{3600}{A} inductive load when a \SI{3}{kV} IGBT static switch
opens. This note quantifies the electrical stress on the crowbar thyristor,
justifies the choice of a \SI{6.5}{kV} symmetric phase-control thyristor (Hitachi
Energy 5STP~26N6500), and explains why the intrinsic self-firing mechanisms of a
thyristor (critical $\mathrm{d}v/\mathrm{d}t$ and forward breakover) must not be
used as they are, but can be turned into a reliable, fully passive trigger by a
floating anode-to-gate network. Design equations and starting values for that
network are given, together with the coordination rules with the MOV clamp of the
main switch and the checks to be made in simulation.
\end{abstract}

% ==========================================================================
\section{Scope and circuit under study}
% ==========================================================================

The circuit is the static-switch test bench of Fig.~\ref{fig:circuit}: a DC source
$V_s = \SI{3000}{V}$ feeds an RL load ($R = \SI{0.8}{\ohm}$, $L = \SI{2}{mH}$)
through an IGBT breaker $S$ (four 5SNA~1000G650300 in parallel). The breaker is
opened under load at $I_0 = \SI{3600}{A}$; a varistor bank across the breaker
clamps its voltage at $V_{\mathrm{MOV}} \approx \SI{5.5}{kV}$. The load energy is
extracted by the crowbar branch in parallel with the load: thyristor $T$, dump
inductor $L_c = \SI{12.5}{\micro H}$ and dump resistor $R_D = \SI{0.42}{\ohm}$,
with an RC snubber $C_{sn} = \SI{2}{\micro F}$, $R_{sn} = \SI{64}{\ohm}$ across the
branch. The bench must perform of the order of $10^6$ opening cycles.

\begin{figure}[htbp]
\centering
\begin{circuitikz}[american, scale=0.95, transform shape]
  % source (positive terminal towards the breaker)
  \draw (0,0) to[battery1, l_=$V_s$] (0,4);
  \node[font=\scriptsize] at (0.35,3.35) {$+$};
  \node[font=\scriptsize] at (0.35,0.65) {$-$};
  % breaker with MOV
  \draw (0,4) -- (1,4) to[nos, l=$S$] (3.5,4) -- (4.5,4);
  \draw (1,4) -- (1,5.5) to[varistor, l=MOV] (3.5,5.5) -- (3.5,4);
  % load
  \draw (4.5,4) -- (5.5,4) to[R, l_=$R$] (5.5,2.5) to[L, l_=$L$] (5.5,0.5) -- (5.5,0) -- (0,0);
  % crowbar branch: thyristor (anode at bottom), Lc, RD
  \draw (5.5,4) -- (8,4);
  \draw (5.5,0) -- (8,0);
  \draw (8,0) to[Ty, n=T, l_=$T$] (8,1.8) to[L, l_=$L_c$] (8,2.9) to[R, l_=$R_D$] (8,4);
  % snubber across the branch
  \draw (8,4) -- (10,4) to[R, l=$R_{sn}$] (10,2.4) to[C, l=$C_{sn}$] (10,0) -- (8,0);
  % gate network box
  \draw (T.G) -- (7.15,0 |- T.G);
  \draw[dashed] (5.95,0.3) rectangle (7.15,1.9);
  \node[align=center,font=\scriptsize] at (6.55,1.1) {gate\\network\\(Fig.~\ref{fig:gate})};
  % current arrow
  \draw[-latex] (4.55,3.75) -- (5.15,3.75) node[midway, above, font=\scriptsize] {$i_L$};
\end{circuitikz}
\caption{Static-switch test bench with dump-resistor crowbar. In service ($S$ closed) the load
voltage is $+V_s$ and the thyristor blocks in the \emph{reverse} direction; forward voltage appears
across $T$ only while $S$ opens.}
\label{fig:circuit}
\end{figure}

Two facts shape the whole analysis:

\begin{enumerate}
\item In service the crowbar thyristor is \emph{reverse} biased at $-V_s = \SI{-3}{kV}$.
      Forward voltage appears only when $S$ opens and the load inductance reverses the load
      voltage, which is exactly when the crowbar must fire.
\item The forward voltage available to the crowbar is bounded by the MOV clamp of the breaker:
      \begin{equation}
        v_{AK,\max} = V_{\mathrm{MOV}} - V_s \approx \SI{5.5}{kV} - \SI{3.0}{kV} = \SI{2.5}{kV}.
        \label{eq:vfw}
      \end{equation}
      Any firing threshold of the crowbar must lie well below this value, otherwise the MOV
      takes the load current and the energy.
\end{enumerate}

% ==========================================================================
\section{Electrical stress on the crowbar thyristor}
% ==========================================================================

\subsection{Conduction: charge, $\int i^2\,\mathrm{d}t$, energy}

Once the load current has commutated into the crowbar branch, it decays through
$R + R_D$ with
\begin{equation}
  \tau = \frac{L}{R + R_D} = \frac{\SI{2}{mH}}{\SI{1.22}{\ohm}} \approx \SI{1.64}{ms}.
\end{equation}
The conduction quantities follow directly:
\begin{align}
  Q &= I_0\,\tau \approx \SI{5.9}{C}, \\
  \int i^2\,\mathrm{d}t &= \frac{I_0^2\,\tau}{2} \approx \SI{10.6e3}{A^2 s}, \\
  E_{L} &= \tfrac12 L I_0^2 \approx \SI{13.0}{kJ},
  \qquad E_{R_D} = E_L\,\frac{R_D}{R+R_D} \approx \SI{4.5}{kJ},
  \qquad E_{R} \approx \SI{8.5}{kJ}.
\end{align}
The thyristor conduction energy, with the on-state model $v_T = V_{T0} + r_T\,i$ of the
5STP~26N6500 ($V_{T0} = \SI{1.12}{V}$, $r_T = \SI{0.29}{m\ohm}$), is
\begin{equation}
  E_T = V_{T0}\,Q + r_T \int i^2\,\mathrm{d}t \approx 6.6 + 3.1 \approx \SI{10}{J}.
\end{equation}
Spread over the silicon of a \SI{100}{mm} wafer (thermal capacity of the order of
\SI{10}{J/K}), this gives $\Delta T_j \lesssim \SI{1}{K}$ per event. The crowbar
thyristor is therefore \emph{not} sized by current or by $\int i^2\mathrm{d}t$; it is sized by
blocking voltage, turn-on $\mathrm{d}i/\mathrm{d}t$ and cycle count. The current reaches the
holding level ($I_H \approx \SI{1}{A}$) after $\tau\ln(I_0/I_H) \approx \SI{13}{ms}$ and the
device turns off naturally, without reverse recovery.

\subsection{Commutation: $\mathrm{d}i/\mathrm{d}t$ and MOV energy}

While $S$ opens, the breaker voltage is held at $V_{\mathrm{MOV}}$ by the varistor bank; the
forward voltage \eqref{eq:vfw} drives the load current into $L_c$:
\begin{align}
  \frac{\mathrm{d}i}{\mathrm{d}t} &= \frac{V_{\mathrm{MOV}} - V_s}{L_c}
      = \frac{\SI{2.5}{kV}}{\SI{12.5}{\micro H}} = \SI{200}{A/\micro s}, \\
  t_{\mathrm{comm}} &= \frac{L_c\,I_0}{V_{\mathrm{MOV}} - V_s} = \SI{18}{\micro s}, \\
  E_{\mathrm{MOV,comm}} &\approx \tfrac12\,V_{\mathrm{MOV}}\,I_0\,t_{\mathrm{comm}} \approx \SI{180}{J}.
\end{align}
The MOV bank thus sees about \SI{180}{J} per opening instead of the $\sim$\SI{16}{kJ} it would
absorb without the crowbar, which is what makes the $10^6$-cycle target credible.
The turn-on $\mathrm{d}i/\mathrm{d}t$ is set by the MOV clamp level and $L_c$, not by the
triggering method: it is the same whether the thyristor fires at \SI{50}{V} or at \SI{2}{kV}.
It must be compared with the \emph{repetitive} critical $\mathrm{d}i/\mathrm{d}t$ of the
datasheet (for large phase-control thyristors of the order of a few hundred \si{A/\micro s});
if the margin is insufficient, increase $L_c$ or lower the MOV clamp.

\subsection{$\mathrm{d}v/\mathrm{d}t$ at breaker opening}

At the instant $S$ turns off, the load current initially flows into the capacitances hanging on
the load node; with the branch snubber alone as an upper bound,
\begin{equation}
  \left.\frac{\mathrm{d}v}{\mathrm{d}t}\right|_{\max} \approx \frac{I_0}{C_{sn}}
  = \frac{\SI{3600}{A}}{\SI{2}{\micro F}} = \SI{1.8}{kV/\micro s},
  \label{eq:dvdt}
\end{equation}
so the thyristor voltage swings from \SI{-3}{kV} through zero into forward bias in about
\SI{2}{\micro s}. The breaker's own RCD snubbers in parallel reduce this figure, but it remains
of the same order as the critical $\mathrm{d}v/\mathrm{d}t$ of a phase-control thyristor
(typically \SIrange{1000}{2000}{V/\micro s}). This is the key observation of
Section~\ref{sec:selffire}: left to itself, the device is likely to self-fire at breaker
opening, in an uncontrolled way.

% ==========================================================================
\section{Device selection}
% ==========================================================================

\subsection{Requirements}

\begin{table}[htbp]
\centering
\caption{Crowbar thyristor requirements and 5STP~26N6500 data.}
\label{tab:req}
\begin{tabular}{@{}llll@{}}
\toprule
Quantity & Requirement & 5STP 26N6500 & Note \\
\midrule
Reverse blocking, continuous DC & \SI{3.0}{kV} & $V_{RRM} = \SI{6.5}{kV}$ & cosmic-ray FIT, see text \\
Forward blocking, transient & $\le$ \SI{2.5}{kV} & $V_{DRM} = \SI{6.5}{kV}$ & bounded by MOV \\
Peak current & \SI{3600}{A} & $I_{TSM} = \SI{65}{kA}$ & \\
$\int i^2\mathrm{d}t$ per event & \SI{10.6e3}{A^2 s} & $\sim\!\SI{2e7}{A^2 s}$ & from $I_{TSM}$, \SI{10}{ms} \\
Turn-on $\mathrm{d}i/\mathrm{d}t$, repetitive & \SI{200}{A/\micro s} & datasheet & check margin \\
$\mathrm{d}v/\mathrm{d}t$ at turn-on & up to \SI{1.8}{kV/\micro s} & \SIrange{1000}{2000}{V/\micro s} & gate network required \\
Cycles & $\sim 10^6$ & --- & $\Delta T_j \lesssim \SI{1}{K}$ \\
Gate trigger current & --- & $I_{GT} \le \SI{0.4}{A}$ & \\
\bottomrule
\end{tabular}
\end{table}

\subsection{Voltage class: why \SI{6.5}{kV}}

The device blocks \SI{3}{kV} DC continuously in the reverse direction for the whole life of the
bench. For high-voltage silicon devices the dominant random failure mechanism under continuous
DC blocking is cosmic-ray induced breakdown, whose rate rises very steeply with the ratio of
applied to rated voltage. A \SI{4.5}{kV} device would be operated at about \SI{67}{\percent} of
its rating, a \SI{6.5}{kV} device at \SI{46}{\percent}. The FIT-versus-$V_{DC}$ curves of the
manufacturer must be consulted for the chosen part, but the \SI{6.5}{kV} class is the natural
choice; it also gives comfortable margin against the transients at breaker closing, when the
device sees a \SI{0}{}$\rightarrow$\SI{-3}{kV} step.

The 5STP~26N6500 (Hitachi Energy, formerly ABB) is a symmetric \SI{6.5}{kV} phase-control
thyristor with interdigitated amplifying gate, $I_{T(RMS)} = \SI{4410}{A}$,
$I_{TSM} = \SI{65}{kA}$ \cite{ds26n6500}. A Simscape model of this device with Lauritzen--Ma
reverse recovery already exists, so the simulation work carries over. Its current margins in this
application are very large (three orders of magnitude on $\int i^2\mathrm{d}t$); a smaller wafer of
the same \SI{6.5}{kV} family would be adequate thermally and would reduce clamping force and cost.
The \SI{100}{mm} device is retained here for the $\mathrm{d}i/\mathrm{d}t$ margin and because it
is already modelled.

\subsection{Why not an asymmetric pulse-power thyristor}

Pulse-power thyristors (e.g.\ the Dynex range based on GTO structures, \SIlist{3300;4500}{V},
highly interdigitated gate, very high $\mathrm{d}i/\mathrm{d}t$) are marketed also for crowbar
duty \cite{dynexpp}. They do not fit here for two reasons: they are asymmetric and do not block
the \SI{-3}{kV} seen in service (a series diode would be needed), and the $\mathrm{d}i/\mathrm{d}t$
of \SI{200}{A/\micro s} imposed by $L_c$ is well within the reach of a phase-control device. A
symmetric PCT is the simpler and more robust choice.

% ==========================================================================
\section{Self-firing mechanisms and why they must not be used as they are}
\label{sec:selffire}
% ==========================================================================

\subsection{$\mathrm{d}v/\mathrm{d}t$ firing}

A rising forward voltage drives a displacement current $C_{J2}\,\mathrm{d}v/\mathrm{d}t$ through
the blocking junction $J_2$. This current flows laterally in the p-base towards the cathode
shorts; where the lateral voltage drop across the p-base under the $n^+$ emitter reaches
$\approx \SI{0.7}{V}$, the emitter injects and the thyristor turns on \cite{baliga}. The
location is wherever the p-base resistance is highest or the shorts are farthest, i.e.\ a random
point of the cathode area, not the gate region. The datasheet ``critical $\mathrm{d}v/\mathrm{d}t$''
is the ramp the device is guaranteed \emph{not} to fire on at $T_{j,\max}$; the actual firing
ramp is unspecified, may be several times higher, and falls with temperature because leakage adds
to the displacement current.

\subsection{Forward breakover}

When the forward voltage exceeds $V_{BO}$ the avalanche multiplication at $J_2$ satisfies the
regenerative condition $\alpha_{npn}\,\alpha_{pnp}\,M \to 1$ and the device turns on
\cite{baliga}. $V_{BO}$ is not a datasheet parameter: it lies at roughly
\numrange{1.1}{1.3}$\times V_{DRM}$, varies with temperature and from device to device, and the
turn-on initiates at the point of maximum field (junction termination or gate periphery). The
undesired ignition of a thyristor occurs at a random, localized spot of the cathode
junction and may destroy the device through the concentrated turn-on losses \cite{bodpatent}.

\subsection{Consequences in this circuit}

\begin{itemize}
\item In both mechanisms the initially conducting area is small and unrelated to the
      amplifying gate; the $\mathrm{d}i/\mathrm{d}t$ that follows (\SI{200}{A/\micro s} here) must
      then be carried by that small area. The datasheet $\mathrm{d}i/\mathrm{d}t$ rating does not
      apply. Over $10^6$ cycles this is not an acceptable operating mode.
\item Static self-firing on $V_{BO}$ is not even reachable: the forward voltage is limited to
      \SI{2.5}{kV} by the MOV \eqref{eq:vfw}, far below the $V_{BO} > \SI{6.5}{kV}$ of the chosen
      device. The MOV would absorb the full \SI{16}{kJ} and the crowbar would never fire.
\item Dynamic self-firing is, on the contrary, \emph{likely} to happen: \eqref{eq:dvdt} is at the
      critical $\mathrm{d}v/\mathrm{d}t$ level. Without a gate network the device will fire, at an
      unpredictable point of the wafer, at an unpredictable forward voltage.
\end{itemize}

The useful conclusion is not to suppress these mechanisms but to \emph{pre-empt} them: the same
anode voltage that would cause an uncontrolled ignition is used to inject a proper gate current,
so that turn-on always starts at the amplifying gate. This is the ``protective firing'' concept of
HVDC valves, where a break-over diode (BOD) between anode and gate converts an overvoltage into a
regular gate pulse \cite{bodpatent,arrillaga}. Here the roles are extended: the dynamic path is
the primary trigger, the static path is the back-up.

% ==========================================================================
\section{Passive self-triggering gate network}
% ==========================================================================

\subsection{Topology}

\begin{figure}[htbp]
\centering
\begin{circuitikz}[american, scale=1.0, transform shape]
  % thyristor vertical, anode on top; gate lead comes out on the right
  \draw (0,6) node[above left]{A} to[Ty, n=T] (0,0) node[below left]{K};
  % anode rail with D1 and node N
  \draw (0,6) -- (1,6) to[D, l=$D_1$] (3,6) -- (8.6,6);
  \node[above] at (3,6) {N};
  \filldraw (3,6) circle (1pt);
  % gate rail
  \draw (T.G) -- (1.5,0 |- T.G) coordinate (G);
  \filldraw (G) circle (1pt);
  \draw (G) -- (3,0 |- G) to[R, l_=$R_{gs}$] (3,4.2) coordinate (G1);
  \filldraw (G1) circle (1pt);
  % bus between the two paths
  \draw (3,4.2) -- (8.6,4.2);
  % dynamic path: Cg with reset resistor RC
  \draw (3,6) to[C, l_=$C_g$] (3,4.2);
  \draw (4.5,6) to[R, l_=$R_C$] (4.5,4.2);
  % static path: BOD / TVS stack and Rs
  \draw (8.6,6) to[zD, l_=BOD] (8.6,5.1) to[R, l_=$R_s$] (8.6,4.2);
  % gate-cathode network: RGK and antiparallel DGK
  \draw (G) to[R, l_=$R_{GK}$] (1.5,0) -- (0,0);
  \draw (1.5,0) -- (2.4,0) to[D, l_=$D_{GK}$] (2.4,0 |- G) -- (G);
  % annotations
  \node[font=\scriptsize, align=left, anchor=west] at (5.0,5.1) {dynamic path\\$i_G = C_g\,\mathrm{d}v/\mathrm{d}t$};
  \node[font=\scriptsize, align=left, anchor=west] at (8.95,5.1) {static path\\fires at $v_{AK} > V_{st}$};
\end{circuitikz}
\caption{Floating anode-to-gate trigger network. $D_1$ isolates the network during reverse
blocking; $C_g$ converts the forward $\mathrm{d}v/\mathrm{d}t$ into gate current; the BOD (or TVS
stack) with $R_s$ fires the gate if the forward voltage rises slowly beyond $V_{st}$;
$R_C$ resets $C_g$; $R_{GK}$ and $D_{GK}$ protect the gate.}
\label{fig:gate}
\end{figure}

All components are referenced to the cathode and float with it: no auxiliary supply, no
fibre-optic link, no isolation transformer. The network works as follows.

\paragraph{$D_1$ -- reverse isolation.} With $v_{AK} < 0$ (all of service life) $D_1$ is reverse
biased and node~N sits at cathode potential: no current reaches the gate. This is essential,
because a positive gate current while the thyristor blocks reverse voltage forward-biases $J_3$
and $J_2$ and is amplified by the pnp transistor gain into an increased reverse leakage
concentrated near the gate; manufacturers exclude this condition. $D_1$ also keeps the
static-path elements from conducting in their forward direction at \SI{-3}{kV}. Its blocking
voltage must cover $V_s$ plus closing transients with margin: a string rated $\ge \SI{6}{kV}$
(e.g.\ 3--4 fast diodes of \SIrange{1.6}{2.2}{kV} with sharing resistors), or a HV diode stack.
Its current is only the gate current, tens of amperes for microseconds.

\paragraph{Dynamic path -- $C_g$.} As soon as $v_{AK}$ crosses zero and rises, $D_1$ conducts and
\begin{equation}
  i_G \approx C_g\,\frac{\mathrm{d}v_{AK}}{\mathrm{d}t}
\end{equation}
flows into the gate through $R_{gs}$ (valid for $R_{gs}C_g$ much shorter than the rise time).
Firing occurs when the injected current exceeds the gate trigger current plus what $R_{GK}$
diverts:
\begin{equation}
  \left.\frac{\mathrm{d}v}{\mathrm{d}t}\right|_{th} = \frac{I_{GT} + V_{GT}/R_{GK}}{C_g}.
  \label{eq:dvth}
\end{equation}
The thyristor therefore turns on at the amplifying gate within a few tens of volts of forward
bias, long before the intrinsic $\mathrm{d}v/\mathrm{d}t$ mechanism could act.

\paragraph{$R_C$ -- reset.} $C_g$ charges to the forward voltage present at firing (up to
\SI{2.5}{kV}); after turn-on $D_1$ blocks and $C_g$ discharges through $R_C$ with
$\tau_{C} = R_C C_g$, which must be much longer than the voltage rise (\si{\micro s}) and much
shorter than the cycle period. $R_C$ dissipates $\tfrac12 C_g v_{AK}^2$ per event.

\paragraph{Static path -- BOD or TVS stack with $R_s$.} If the forward voltage rises too slowly for
the dynamic path (abnormal breaker opening, fault), the BOD conducts once
$v_{AK} > V_{st}$ and drives the gate through $R_s$, which limits the current to
$(v_{AK} - V_{st})/R_s$. Since no forward voltage exists in service, $V_{st}$ can be low; it must
only exceed ringing and noise and remain far below \eqref{eq:vfw}. Genuine BODs are available from
HVDC gate-unit suppliers; at a threshold of a few hundred volts a stack of high-energy TVS diodes
(5KP/15KP class) does the same job.

\paragraph{$R_{GK}$, $D_{GK}$.} $R_{GK}$ gives noise immunity and sets the threshold
\eqref{eq:dvth}; $D_{GK}$ (antiparallel) limits negative gate voltage to well below $V_{RGM}$,
e.g.\ against the small displacement current through the $D_1$ junction capacitance at breaker
closing.

\subsection{Starting values}

\begin{table}[htbp]
\centering
\caption{Starting values for the gate network (5STP~26N6500, $I_{GT} \le \SI{0.4}{A}$,
$V_{GT} \approx \SI{3}{V}$ to be confirmed on the datasheet).}
\label{tab:gate}
\begin{tabular}{@{}llp{7.2cm}@{}}
\toprule
Element & Value & Design check \\
\midrule
$C_g$ & \SI{4.7}{nF}, $\ge \SI{4}{kV}$ DC & $i_G = \SI{2.4}{A}$ at \SI{500}{V/\micro s}, \SI{8.5}{A} at \SI{1.8}{kV/\micro s}; compare peak with $I_{FGM}$ \\
$R_{GK}$ & \SI{10}{\ohm} & threshold \eqref{eq:dvth}: $(0.4 + 0.3)\,\si{A}/\SI{4.7}{nF} \approx \SI{150}{V/\micro s}$ \\
$R_{gs}$ & \SIrange{10}{15}{\ohm} & $R_{gs}C_g \approx \SI{70}{ns} \ll$ rise time; damps the $C_g$--lead-inductance loop, $2\sqrt{L_{lead}/C_g} \approx \SI{13}{\ohm}$ for $L_{lead} = \SI{0.2}{\micro H}$ \\
$R_C$ & \SI{22}{k\ohm}, HV & $\tau_C \approx \SI{100}{\micro s}$; $E = \tfrac12 C_g (\SI{2.5}{kV})^2 \approx \SI{15}{mJ}$ per event \\
BOD / TVS & $V_{st} \approx \SIrange{400}{500}{V}$ & e.g.\ 3 $\times$ \SIrange{150}{180}{V} TVS 5KP/15KP; $V_{st} \ll \SI{2.5}{kV}$ \\
$R_s$ & \SI{200}{\ohm}, pulse rated & $i_{G,\max} = (2500 - 500)/200 = \SI{10}{A}$; $\sim\!\SI{15}{mJ}$ in the stack per event (worst case) \\
$D_1$ & string $\ge \SI{6}{kV}$, fast & blocks $-V_s$ in service; sharing R/C across the string \\
$D_{GK}$ & fast or Schottky, $\ge \SI{15}{A}$ pulse & clamps $v_{GK} < 0$ \\
\bottomrule
\end{tabular}
\end{table}

\subsection{Optional active trigger}

If the breaker controller can issue a gate command, it should be used as the primary trigger
(pulse transformer or fibre-optic gate unit) with the passive network of Fig.~\ref{fig:gate} left
in place as an independent back-up -- the same philosophy as HVDC valves. The command must be
issued together with, not before, the IGBT turn-off: a positive gate pulse applied while the
thyristor still blocks \SI{-3}{kV} produces the reverse-leakage amplification described above.

% ==========================================================================
\section{Coordination with the MOV clamp}
% ==========================================================================

\begin{enumerate}
\item \textbf{Firing level:} $V_{st} \ll V_{\mathrm{MOV}} - V_s$. With $V_{st} \approx \SI{0.5}{kV}$
      and \SI{2.5}{kV} available the margin is a factor five; the dynamic path fires even
      earlier.
\item \textbf{$\mathrm{d}i/\mathrm{d}t$:} $(V_{\mathrm{MOV}} - V_s)/L_c = \SI{200}{A/\micro s}$,
      independent of the trigger. To be checked against the repetitive datasheet value; the two
      knobs are $L_c$ and the MOV clamp level.
\item \textbf{MOV duty:} about \SI{180}{J} per opening during the \SI{18}{\micro s} commutation
      interval; the varistor bank becomes a back-up clamp with negligible ageing.
\item \textbf{Branch snubber:} at firing $C_{sn}$ discharges through $R_{sn}$ into the thyristor,
      $v_{AK}/R_{sn} \le \SI{40}{A}$ and $\tfrac12 C_{sn} v_{AK}^2 \le \SI{6}{J}$ in $R_{sn}$:
      harmless.
\end{enumerate}

% ==========================================================================
\section{Verification in simulation}
% ==========================================================================

With the existing Simscape model of the 5STP~26N6500 and the gate network added:

\begin{itemize}
\item Plot $v_{AK}$ (device terminals, not branch voltage), $i_G$ and $i_A$ at breaker opening.
      Confirm that $i_G$ exceeds $I_{GT}$ while $v_{AK}$ is still low, and record the forward voltage
      at which anode current starts.
\item Extract $\mathrm{d}i_A/\mathrm{d}t$ during the first \SI{20}{\micro s} and compare with the
      repetitive rating.
\item Sweep the effective capacitance at the load node (branch snubber plus breaker snubbers) to
      cover the range of $\mathrm{d}v/\mathrm{d}t$ from \SI{200}{V/\micro s} to \SI{2}{kV/\micro s};
      the static path should take over below the threshold \eqref{eq:dvth}.
\item Breaker closing: $v_{AK}$ steps from \SI{0}{} to \SI{-3}{kV}; check that the gate stays
      quiet ($v_{GK}$ within $\pm$ a few volts) and that $D_1$ sees no overvoltage.
\item Abnormal case: breaker re-closing while the crowbar still conducts. The source then forces
      the crowbar current to zero through $L_c$ at $V_s/L_c = \SI{240}{A/\micro s}$ with reverse
      recovery; check the reverse voltage overshoot across $T$ and the branch snubber. The
      turn-off time $t_q$ is irrelevant in normal operation, since forward voltage returns only at
      the next cycle.
\end{itemize}

% ==========================================================================
\section{Conclusions}
% ==========================================================================

The crowbar thyristor of the \SI{3}{kV} / \SI{3600}{A} bench dissipates about \SI{10}{J} per
event and is sized by voltage, $\mathrm{d}i/\mathrm{d}t$ and cycle count. A symmetric
\SI{6.5}{kV} phase-control thyristor such as the 5STP~26N6500 is the appropriate choice: it
blocks the \SI{3}{kV} reverse DC of service life with adequate cosmic-ray margin, and its
$\mathrm{d}i/\mathrm{d}t$ capability matches the \SI{200}{A/\micro s} imposed by $L_c$ and the
MOV clamp. The intrinsic self-firing mechanisms must not be relied upon: forward breakover is
unreachable below the MOV clamp, and $\mathrm{d}v/\mathrm{d}t$ firing is likely but uncontrolled.
A floating anode-to-gate network -- blocking diode, coupling capacitor, BOD/TVS stack -- turns
both into a regular gate-triggered turn-on with no auxiliary supply, and can be combined with
an active command as primary trigger. Firing level and $\mathrm{d}i/\mathrm{d}t$ are coordinated
with the MOV clamp through $V_{st} \ll V_{\mathrm{MOV}} - V_s$ and
$(V_{\mathrm{MOV}} - V_s)/L_c$.

% ==========================================================================
\begin{thebibliography}{9}

\bibitem{ds26n6500}
Hitachi Energy (formerly ABB Semiconductors), ``Phase Control Thyristor 5STP~26N6500,'' data sheet.

\bibitem{baliga}
B.~J. Baliga, \emph{Fundamentals of Power Semiconductor Devices}. New York: Springer, 2008,
ch.~8 (Thyristors).

\bibitem{bodpatent}
``Protective circuitry for a thyristor,'' U.S. Patent 4\,651\,251 (break-over diode protective firing, Nelson River HVDC precedent).

\bibitem{arrillaga}
J.~Arrillaga, \emph{High Voltage Direct Current Transmission}, 2nd ed. London: IEE, 1998
(thyristor valves, protective firing).

\bibitem{dynexpp}
Dynex Semiconductor, ``Pulsed power thyristors,'' product information, accessed Sept.~2026.

\bibitem{abbsnubber}
ABB Switzerland Ltd, Semiconductors, ``Design of RC snubbers for phase control applications,''
Application Note 5SYA~2020.

\end{thebibliography}

\end{document}
